\chapter{Chatbots, Models and Agents}
\label{ch:landscape}

Three words get used interchangeably and mean quite different things. Getting
them apart is the first useful step.

\section{The model}

A large language model is a function. It takes a sequence of tokens and returns
a probability distribution over the next token. That is the whole interface.
Everything else --- conversation, tool use, reasoning, refusal --- is built on
top of that one operation, repeated.

The model itself is stateless. It has no memory of your previous question, no
file system, no clock, and no ability to act. Each call is independent. If a
system appears to remember something, that is because the text was fed back in.

\begin{keypoint}
The model is a calculator, not a colleague. It has no persistent state. Anything
that looks like continuity is machinery around the model, not in it.
\end{keypoint}

\subsection{What ``training'' produced}

Pretraining fits the next-token function on a very large corpus. This is where
the model acquires syntax, factual associations, code idioms, and the statistical
shape of academic prose. A second stage --- instruction tuning and preference
optimisation --- reshapes the outputs to be helpful, to follow instructions, and
to decline certain requests.

That second stage matters more than it looks. It is where sycophancy comes from
(\Cref{ch:powers}), where refusal behaviour is installed, and where the model
learns to guess confidently rather than say it does not know
\parencite{kalai2025hallucinate}.

\subsection{Emergence, and whether it is real}

Claims that language models display \term{emergent} capabilities --- abilities
that appear abruptly at scale rather than improving smoothly --- have been
central to the field's self-description. \textcite{krakauer2026emergence} examine
these claims from a complex systems perspective and are careful about them:
much apparent emergence is an artefact of discontinuous evaluation metrics, and
the deeper question of whether LLMs possess emergent \emph{intelligence} rather
than emergent \emph{capability} remains open.

\textcite{mitchell2026jagged} names the practical version of this problem:
\term{jagged intelligence}. Performance is not a single number. A system that
handles a research-level reasoning task can fail at counting, at a spatial
relation, or at a variant of the same question with different surface form. The
jaggedness is not noise around a trend; it is the shape of the capability
itself. This is why benchmark scores translate so poorly into expectations about
your own task.

There is also a formal strand arguing that the architecture is more legible than
its reputation suggests. \textcite{coppola2026transformers} proves that a sigmoid
transformer implements weighted loopy belief propagation on an implicit factor
graph --- one layer is one round of belief propagation. Whether this reframing
changes anything in practice is unsettled, but it does undercut the claim that
transformers cannot be understood.

\section{The chatbot}

A chatbot is a model plus a conversation buffer plus a web interface. The buffer
is the important part: it concatenates the history and resubmits it on every
turn. This is why long conversations get expensive, why they eventually forget
the beginning, and why starting fresh often fixes a stuck session.

Chatbots are good for: drafting, explaining, translating, rephrasing,
brainstorming, reading a paper you would otherwise skim. They are bad for
anything requiring verified fact, current information, or action in the world.

\section{The agent and the harness}

An \term{agent} is a model given tools and a loop. An \term{agent harness} is the
orchestration system that manages that loop: it decides what the model sees,
executes the tool calls the model requests, feeds results back, enforces
permissions, and decides when to stop.

The core is unremarkable. \textcite{liu2026claudecode} analysed the architecture
of a production coding agent and found the centre of the system to be a simple
while-loop: call the model, run the tools it asks for, repeat. Most of the code
is not in the loop. It is in the systems around it --- the permission layer, the
context manager, the memory, the error recovery.

\begin{lstlisting}[style=shell,caption={The whole idea, in pseudocode}]
while not done:
    response = model(context)
    if response.wants_tool:
        result = run_tool(response.tool, response.args)   # <- permissions here
        context += result
    else:
        done = True
\end{lstlisting}

That single \code{run\_tool} line is where the security risk lives. It is the
point at which a text generator becomes something that can delete your files.

\subsection{Why the harness matters more than the model}

Two harnesses running the same model behave very differently. The differences
come from:

\begin{itemize}
  \item \textbf{Context assembly} --- what goes into the prompt, in what order,
    and what gets dropped when space runs out.
  \item \textbf{Tool design} --- a good file-edit tool that fails loudly beats a
    clever one that fails silently.
  \item \textbf{Permission model} --- what runs without asking, what prompts,
    what is refused outright.
  \item \textbf{Memory} --- what persists across sessions, and in what form.
\end{itemize}

\textcite{du2026memory} surveys agent memory as a \term{write--manage--read} loop
coupled to perception and action, spanning context-resident compression,
retrieval-augmented stores, reflective self-improvement and hierarchical virtual
context. The taxonomy is useful, but the practical finding for our purposes is
simpler: the memory that works most reliably today is a plain file on disk that
the agent reads at the start of every session.

\subsection{Files are all you need}

\textcite{piskala2026files} traces this explicitly to Unix. The 1970s abstraction
--- \emph{everything is a file} \parencite{ritchie1974unix} --- collapsed
heterogeneous resources into one uniform interface, and it survived fifty years
of hardware change. Agents face the analogous problem: databases, APIs, cloud
infrastructure, memories, human input, each with its own protocol. The same
solution works. Designing an agent around file-like abstractions makes it
inspectable, versionable, greppable and diffable.

This has a direct consequence for how you set up a project. A \code{README.md}
that explains the repository and an \code{AGENTS.md} (or \code{CLAUDE.md}) that
states the conventions will do more for agent performance than most prompt
engineering. They are plain text, they live in version control, and you can read
them yourself.

\begin{keypoint}
\code{AGENTS.md} is not a wish list. It is a behavioural contract that closes
specific failure modes you have actually observed. Write it after something goes
wrong, not before.
\end{keypoint}

\section{The ecosystem, briefly}

The tooling layer is crowded and turns over fast. As of writing, the
terminal-first agent harnesses in common use include Claude Code, OpenCode,
Codex, Gemini CLI, Copilot CLI and Aider; IDE-embedded ones include Cline,
Continue and Roo Code. Domain-specific agents exist for the life sciences:
Biomni \parencite{huang2025biomni}, PantheonOS \parencite{xu2026pantheonos},
OmicClaw \parencite{zeng2026omicclaw}, and others appearing monthly.

Names will change. The design --- model, loop, tools, permissions, files --- has
been stable since the pattern emerged, and is the thing worth learning. Pick one
that runs in your terminal, works with the models you can afford, and shows you
what it is about to do. That is most of the decision.

\section{Where agents are going}

Two directions are worth watching because they bear on research practice
directly.

The first is autonomous research. \textcite{lu2026aiscientist} report a system
that navigates the research process end to end --- idea generation, experiment,
writing --- and \textcite{feng2026aletheia} describe a mathematics research agent
that has produced genuine results. \textcite{song2026paperorchestra} automate
manuscript production from raw materials. \textcite{qu2026coral} push toward
multi-agent evolution for open-ended discovery. Whether this scales to field
ecology, where the bottleneck is fieldwork rather than analysis, is a separate
question.

The second is the benchmark reality check. \textcite{guo2026promptbio} evaluated
bioinformatics agents across 244 expert-curated tasks and found performance
declining sharply with task difficulty. The honest summary is that agents are
useful for scoped, well-specified analytical work and unreliable for anything
requiring sustained judgement. That gap is where a domain expert remains
necessary, and it is not obviously closing.
